\chapter*{Conclusion}
\addcontentsline{toc}{chapter}{Conclusion}
\markboth{Conclusion}{Conclusion}\label{ch:conclusion}

This book set out to ask one question: how far can bare-metal Ada go on its own? On
the evidence of the preceding chapters, the answer is most of the way to a finished
product. A 2nd-stage bootloader brings the chip up, a from-scratch runtime starts
both cores, and roughly fifty hand-written drivers reach the peripherals. On top of
those sit an ext4 filesystem checked against the Linux kernel's own tools, a TCP/IP
stack over the W5500, and a pure-Ada TLS~1.3 client that fetches a live forecast over
real HTTPS, all with no ESP-IDF and no C cryptography library underneath. The pieces
a microcontroller project usually borrows from a vendor SDK were, here, written in
Ada and read end to end.

\section*{What the language bought}

The reason to do that was never novelty. It was what the language gives back.
Range-checked types reject a wrong-width register write or an out-of-range pin before
the program runs. A private type guards every hardware resource, so a transfer cannot
reach the registers without proving ownership (Chapter~\ref{ch:packages}). The Jorvik
profile keeps the task set fixed and the whole program statically analysable, which
is why the worst-case depth of every stack, and every byte of flash and RAM, can be
read off before the board powers on (Chapter~\ref{ch:static-mem}). The hard logic,
the allocator, the filesystem, the interpreter, runs on the host against a trusted
oracle because it is written target-free (Chapter~\ref{ch:testing}). The recurring
result was fewer surprises at run time, paid for once at compile time.

\section*{What it cost}

It would be dishonest to end on only the wins. The cost of owning the whole stack is
owning the whole stack. Everything the SDK hands you for free, the boot sequence, the
drivers, the cryptography, is yours to write and to debug, and some of that debugging
was real. The PSRAM would not run at full speed until its data-in timing was measured
on the board itself (Chapter~\ref{ch:psram}). The Xtensa register windows,
mismanaged, corrupted the call chain in ways that surfaced chapters apart
(Chapter~\ref{ch:context-switch}). A cellular modem that refused to answer turned out
to be a level shifter on the board, not a fault in any line of Ada, and it took a
single-pad loopback to prove it (Chapter~\ref{ch:debug}). Bare metal has sharp edges
of its own, the elaboration order and the constraints of the leanest profile, that a
hosted runtime hides. None of this argues against the approach. It argues for going
in with eyes open.

\section*{When to reach for this}

So the honest recommendation is conditional. Reach for bare-metal Ada when the system
has to be trusted: when it is safety-relevant, when it must run for years unattended,
when a regulator or your own conscience will ask you to account for every path, when
the cost of a field failure dwarfs the cost of building the foundation yourself. For
a weekend prototype, or a product whose value is in shipping this month, the vendor
SDK is the faster road and there is no shame in taking it. This book is for the other
case, where bounded, auditable, and provable are worth more than fast.

\section*{Where to go from here}

The book is a companion to a living code base, not a closed artifact. Every concept here is anchored in code you can clone, build, and run with
one command, and the examples are the place to start changing things. Read Further
Reading (the references that close the book) for the standards and projects this work
stands on, pick the example nearest to what you want to build, and make it do
something new. The proof that bare-metal Ada goes far enough is not in these pages. It
is on the board.
